\usepackage{float}
\usepackage{fancyhdr}
\usepackage{setspace}
\usepackage[titles]{tocloft}
\usepackage{xcolor}
\usepackage{etoolbox}

% ==========================================
% FONTS EMBEDDING CONFIGURATION
% ==========================================
\usepackage{luacode}
\begin{luacode}
  local function embedfull(tfmdata)
    tfmdata.embedding = "full"
  end
  luatexbase.add_to_callback("luaotfload.patch_font", embedfull, "embedfull")
\end{luacode}

% ==========================================
% LINE SPACING CONFIGURATION (APA 7th)
% ==========================================
% Main text: double spacing (200%)
\doublespacing

% Mono text: 1.2x spacing (120%)
\AtBeginEnvironment{lstlisting}{\setstretch{1.2}}
\AtBeginEnvironment{verbatim}{\setstretch{1.2}}
\AtBeginEnvironment{Shaded}{\setstretch{1.2}}

% ==========================================
% HYPERLINK CONFIGURATION
% ==========================================
% Black links without underline (APA compliant)
\usepackage[
  hidelinks,
  breaklinks=true,
  colorlinks=false,
  bookmarks=true,
  bookmarksopen=false,
  linktoc=all
]{hyperref}

% Override link colors to black (redundant but ensures compliance)
\hypersetup{
  linkcolor=black,
  citecolor=black,
  urlcolor=black,
  filecolor=black
}

% ==========================================
% HEADER AND FOOTER CONFIGURATION
% ==========================================
% Setup fancyhdr for running header with short title and page number
\pagestyle{fancy}
\fancyhf{} % Clear default headers/footers

% Right side: page number
\fancyhead[R]{\thepage}
% Left side: see before-body.tex for short title in the page header


% Set header height and line
\setlength{\headheight}{0.5in}
\renewcommand{\headrulewidth}{0.4pt}
\renewcommand{\footrulewidth}{0pt}

% ==========================================
% FIRST PAGE (TITLE PAGE) - NO PAGE NUMBER
% ==========================================
\fancypagestyle{empty}{%
  \fancyhf{}%
  \renewcommand{\headrulewidth}{0pt}%
  \renewcommand{\footrulewidth}{0pt}%
}

% Apply 'empty' style to first page automatically
\AtBeginDocument{%
  \thispagestyle{empty}%
}

% ==========================================
% TABLE OF CONTENTS FORMATTING
% ==========================================
\renewcommand{\cftsecleader}{\cftdotfill{\cftdotsep}}
\renewcommand{\cftsecfont}{\normalfont}
\renewcommand{\cftsecpagefont}{\normalfont}
\renewcommand{\cftsubsecfont}{\normalfont}
\renewcommand{\cftsubsecpagefont}{\normalfont}

% ==========================================
% FIGURE CAPTIONS (APA 7th STYLE)
% ==========================================
\usepackage[
  labelsep=newline,
  justification=justified,
  singlelinecheck=false,
  font=small,
  labelfont=bf,
  textfont=it,
  format=plain
]{caption}

% Create custom format for figures with vertical space between label and title
% This adds 0.2em of space between "Figura 1" and the title text
\DeclareCaptionFormat{apa7figure}{#1\\[0.8em]#3}

% APA 7th Figure Format:
% Figura 1 (bold)
% [vertical space]
% Figure title in italics
% [image]
% Nota. Optional note text (within same figure environment)
\captionsetup[figure]{
  name=Figura,
  labelsep=none,
  justification=justified,
  singlelinecheck=false,
  labelfont=bf,
  textfont=it,
  format=apa7figure,
  position=top
}

% Spanish language for figures
\renewcommand{\figurename}{Figura}
\renewcommand{\tablename}{Tabla}

% ==========================================
% TABLE CAPTIONS (APA 7th STYLE)
% ==========================================
% APA 7th: Table title in italics, on separate line from label, above the table
\captionsetup[table]{
  name=Tabla,
  position=top,
  labelsep=newline,
  justification=justified,
  singlelinecheck=false,
  labelfont=bf,
  textfont=it,
  format=plain
}

% ==========================================
% FIGURE ENVIRONMENT CUSTOMIZATION (APA 7th)
% ==========================================
% Customize spacing around figures to comply with APA 7th edition
\setlength{\abovecaptionskip}{1.0em}
\setlength{\belowcaptionskip}{0.5em}

% Note: Figure positioning (H - here) is controlled by the 'fig-pos' option 
% in _extension.yml configuration, which applies to all Quarto-generated figures.

% Custom command to add figure/table notes within the environment (APA 7th)
% This keeps the note together with the figure/table, preventing page breaks
% Usage: After \includegraphics, use: \figurenote{Note text}
\newcommand{\figurenote}[1]{%
  \vspace{6pt}%
  \par\raggedright\small\textit{Nota.~}#1%
}

% Alternative: For tables, use the same command
\newcommand{\tablenote}[1]{%
  \vspace{6pt}%
  \par\raggedright\small\textit{Nota.~}#1%
}

% ==========================================
% SOURCE CODE LISTINGS CONFIGURATION
% ==========================================
\usepackage{listings}
\usepackage{xcolor}

\definecolor{codebackground}{rgb}{0.95,0.95,0.95}
\definecolor{codenumber}{rgb}{0.5,0.5,0.5}
\definecolor{codestring}{rgb}{0.6,0.1,0.1}
\definecolor{codecomment}{rgb}{0.4,0.4,0.4}
\definecolor{codekeyword}{rgb}{0,0,0.8}

\lstset{
  basicstyle=\ttfamily\small,
  backgroundcolor=\color{codebackground},
  breaklines=true,
  breakatwhitespace=true,
  captionpos=b,
  frame=single,
  frameround=fttt,
  framesep=8pt,
  numbers=left,
  numberstyle=\tiny\color{codenumber},
  numbersep=5pt,
  showstringspaces=false,
  stringstyle=\color{codestring},
  commentstyle=\color{codecomment},
  keywordstyle=\color{codekeyword}\bfseries,
  columns=flexible,
  keepspaces=true,
  aboveskip=1em,
  belowskip=1em,
  abovecaptionskip=0.5em,
  belowcaptionskip=0.5em
}

% Format: "Listado 1. Title text"
\renewcommand{\lstlistingname}{Listado}
\renewcommand{\lstlistlistingname}{Lista de Listados de Código}

\captionsetup[lstlisting]{
  labelsep=period,
  justification=justified,
  singlelinecheck=false,
  labelfont=bf,
  textfont=normalfont,
  font=small,
  skip=0.5em
}

% ==========================================
% EQUATIONS CONFIGURATION
% ==========================================
% amsmath: Professional equation typesetting and multi-line equation support
% amssymb: Extended mathematical symbols and operators
% amsthm: Theorem and proof environments
\usepackage{amsmath}
\usepackage{amssymb}
\usepackage{amsthm}

% Number equations by section (APA style)
% Example: Section 1 equations are numbered 1.1, 1.2, 1.3, etc.
\numberwithin{equation}{section}

% ==========================================
% LIST OF EQUATIONS (OPTIONAL)
% ==========================================
% Controls whether a "List of Equations" appears in the front matter.
% This is controlled by the 'loe' setting in your Quarto format configuration.
% By default, loe: true (list is created). Set loe: false to disable.
%
% USAGE INSTRUCTIONS FOR USERS:
% =============================
% To add an equation to the "List of Equations", use the \myequations{} command.
% This command should be used AFTER the equation environment.
%
% Example 1: Simple equation with description
% ============================================
% \begin{equation}
%   E = mc^2
% \end{equation}
% \myequations{Einstein's mass-energy equivalence}
%
% Example 2: Equation with label (for cross-references)
% ======================================================
% \begin{equation}
%   F = ma \label{eq:newtons-second-law}
% \end{equation}
% \myequations{Newton's Second Law of Motion}
%
% Later in text: Refer to it with \ref{eq:newtons-second-law}
%
% Example 3: Multi-line equations
% ===============================
% \begin{align}
%   y &= mx + b \\
%   0 &= mx + b - y \label{eq:line-form}
% \end{align}
% \myequations{Linear equation in slope-intercept form}
%
% PROFESSIONAL GUIDELINES:
% ========================
% - Only list IMPORTANT or KEY equations in the document
% - All labeled equations can be referenced using \ref{label}
% - The List of Equations only includes entries you explicitly add with \myequations{}
% - This follows APA 7th edition best practices (not all equations are listed)
% - Ensure descriptions are concise and meaningful (max 1-2 lines)

% Setup the list command using tocloft package
\newcommand{\listequationsname}{Lista de Ecuaciones}
\newlistof{myequations}{loe}{\listequationsname}

% Define command for users to call after an equation
% Syntax: \myequations{Description of the equation}
\newcommand{\myequations}[1]{%
  \addcontentsline{loe}{myequations}{\protect\numberline{\theequation}#1}\par}

% Formatting for the list of equations entries
\setlength{\cftmyequationsindent}{0em}       % No indentation for equation entries
\setlength{\cftmyequationsnumwidth}{3em}     % Space for equation numbering

% ==========================================
% SECTION HEADING FORMATTING
% ==========================================
% APA 7th: Left-aligned, bold, centered (for Level 1), etc.
\usepackage{titlesec}

% Level 1: Centered, Bold, Capitalized
\titleformat{\section}
  {\normalfont\normalsize\bfseries\centering}
  {\thesection.~}
  {0em}
  {}[]

% Level 2: Left-aligned, Bold, Capitalized
\titleformat{\subsection}
  {\normalfont\normalsize\bfseries}
  {\thesubsection.~}
  {0em}
  {}[]

% Level 3: Left-aligned, Bold Italic
\titleformat{\subsubsection}
  {\normalfont\normalsize\bfseries\itshape}
  {\thesubsubsection.~}
  {0em}
  {}[]

% Spacing after section headings
\titlespacing*{\section}{0pt}{*3}{*1}
\titlespacing*{\subsection}{0pt}{*3}{*1}
\titlespacing*{\subsubsection}{0pt}{*3}{*1}

% ==========================================
% BLOCKQUOTE FORMATTING
% ==========================================
% Note: Blockquote formatting is handled by Pandoc's built-in quoting package
% No additional customization needed for APA 7th compliance

% ==========================================
% FOOTNOTE/ENDNOTE FORMATTING
% ==========================================
% Single spacing for footnotes
\let\oldfootnote\footnote
\renewcommand{\footnote}[1]{%
  \oldfootnote{\setstretch{1.0}\small #1}%
}
